% ============================================================================
% 4.5 综合到 FPGA
% ============================================================================
\subsection{从 Verilog 到真实硬件的流程}

写好的 Verilog 代码通过下面几个步骤变成可以烧到 FPGA 里的文件：

\begin{enumerate}
    \item \textbf{综合（Synthesis）}：把 Verilog 代码翻译成一个由逻辑门、
        触发器、查找表（LUT）组成的网表
    \item \textbf{布局（Placement）}：决定每个逻辑门放在 FPGA 的哪个位置
    \item \textbf{布线（Routing）}：在 FPGA 的可编程互连资源里
        把需要通信的逻辑门连接起来
    \item \textbf{时序分析（Timing Analysis）}：验证设计能在目标时钟频率下
        正确工作——每条路径的传播延迟必须小于一个时钟周期
    \item \textbf{生成比特流（Bitstream Generation）}：把最终配置信息
        打包成 FPGA 芯片可以理解的二进制文件
\end{enumerate}

你手工用的工具通常是 Vivado（Xilinx 系列 FPGA）、Quartus（Intel/Altera）
或者开源工具链（Yosys + nextpnr + icestorm，针对 Lattice FPGA）。

\subsection{u\_lite 程序如何加载到 FPGA 的指令存储器}

在设计一个自制 CPU 时，指令内存有两种主要的初始化方式：

\paragraph{方式 1：利用 Block RAM 初始化}
在综合时，把你汇编产出的机器码写成一个初始化文件（.coe 或 .mem），
让综合器在编译时把这些值写进 FPGA 的 BRAM 里。

\begin{lstlisting}[style=verilog,caption=BRAM 里嵌入机器码的模板]
// 这个模块在综合时被工具识别为 "Single Port RAM with initial contents"
reg [31:0] instr_mem [0:1023];  // 1024 条 32 位指令 = 4 KB

initial begin
    $readmemh("code.mem", instr_mem);  // 综合时被解析
end

// 每个时钟沿读出一条指令
always @(posedge clk)
begin
    instr_out <= instr_mem[pc[11:2]];
end
\end{lstlisting}

在编译（汇编）层面，你把 \texttt{start.s} 这类文件用
\texttt{clang --target=aarch64-none-elf -c start.s -o start.o} 产出 .o，
再用 \texttt{llvm-objcopy -O binary} 产出纯二进制，最后
写一段脚本把它转成一行一个十六进制数的 .mem 文件。

\paragraph{方式 2：通过 UART bootloader}
更灵活的做法：在 FPGA 设计里放一个小的 bootloader（固定在 BRAM 里），
它在上电后立即启动，通过 UART 接收外部发来的程序二进制，
把它写入一块更大的 RAM/ROM，然后跳过去执行。

这对应 QEMU 章节里你所看到的思路：程序可以在运行时被写入硬件。
UART 发送器（uart\_tx）把程序字节送来，UART 接收器（uart\_rx）
把它们读出来，写进指令存储器——你就不需要重新综合了。

\subsection{和 QEMU 模拟的对应关系}

现在把第 3 章的 QEMU 裸机程序和本章的 Verilog UART 做个对比：

\begin{center}
\begin{tabular}{l l l}
\hline
层面 & QEMU 里 & FPGA 里对应的 \\
\hline
UART 设备 & 虚拟 PL011 & 你写的 UART 模块 \\
地址 & 内存路由到 PL011 & 地址解码器路由到 UART \\
CPU 取指 & TCG 翻译执行 ARM & 你未来写的 CPU 核 \\
DRAM & QEMU 模拟的内存 & FPGA 片上 BRAM 或 DDR \\
\hline
\end{tabular}
\end{center}

你现在写的 UART 模块，加上一个简单的 CPU，就是一个最小的 SoC。
把之前写的 QEMU 汇编程序（读取输入）改写成这个自制 CPU 可以
执行的指令格式后，就能在 FPGA 里跑起来。

\subsection{为什么要做这种事}

从学习角度，写过 UART 之后再回头看 QEMU 里的那段汇编：

\begin{lstlisting}[style=asm]
wait_uart:
    ldr  w3, [x0, #0x18]
    tbnz w3, #5, wait_uart
    str  w2, [x0]
\end{lstlisting}

你会立刻理解它的意思——这是在读取 FLAG 寄存器、检查发送满、
然后写数据寄存器。这些动作不是某个抽象的"系统调用"，
而是真正把字节送到一个你自己用 Verilog 写过的硬件模块里。
你知道电路里发生了什么，就知道软件层为什么要这么写。

反过来说，真正写过软件之后再去写硬件，你也清楚 UART 的数据寄存器
和状态寄存器为什么要分开成两个地址——因为软件会分开读它们。
两者是一体两面。
